\bookchapter{第十五章：電池科學與續航最佳化}

\chapterintro{電池就是你手機的“能量心臟”。這章帶你搞懂電池的小脾氣，學會科學的充電和續航技巧，徹底告別電量焦慮。}

\booksection{一、電池是怎麼回事？迴圈與健康度}

手機電池（鋰離子或鋰聚合物）像個能量海綿：充電就是吸水（存電），放電就是擠水（用電）。有幾個特點值得記住：

\begin{itemize}
  \item \textbf{隨用隨充}：和舊電池不同，它沒有記憶效應，不用非得用完再充。
  \item \textbf{怕熱}：高溫是電池頭號殺手。
  \item \textbf{會老化}：用久了，最大容量就會下降，這是正常現象。
\end{itemize}

\booksubsection{核心指標：迴圈次數與健康度}

\begin{itemize}
  \item 一次迴圈不是充一次電。比如今天用60\%充滿，明天用40\%充滿，總消耗100\%電量才算一次迴圈。
  \item 電池健康度就是當前最大容量佔全新時的百分比。低於80\%時，你會明顯感覺續航變短。
  \item \textbf{怎麼看？}iPhone：設定 > 電池 > 電池健康與充電；安卓：一般去設定裡搜電池健康。
\end{itemize}

\booksection{二、充電玄學？快充、無線充與黃金法則}

\begin{itemize}
  \item \textbf{快充傷電池？}用原裝或認證的快充配件，傷害很小。手機和充電頭有溫控，過熱會自動降功率。
  \item \textbf{無線充電}：方便但效率較低，發熱更明顯。適合整夜充電。
  \item \textbf{黃金養生法則}：儘量讓電量保持在20\%-80\%之間，別老讓它餓著（0\%）或撐著（100\%）。
\end{itemize}

\begin{tcbinfo}{iPhone"80\%充電限制"功能——把黃金法則變成系統設定}
iPhone 15/16系列（iOS 17以上）已經把“20\%-80\%法則”做成了系統功能。開啟後，手機充到80\%會自動停止充電，從而大幅延緩電池老化。設定路徑：

\texttt{設定 > 電池 > 電池健康與充電 > 充電最佳化 > 80\%限制}

三個選項的含義：(1)\textbf{不限}——充到100\%；(2)\textbf{最佳化電池充電}——學習你的作息，只在凌晨最後一段延遲到100\%；(3)\textbf{80\%限制}——直接停在80\%，最護電池。長期在家/辦公室有充電條件的人，強烈建議選80\%限制；經常出差、需要滿電出門的人，可以選最佳化電池充電。

\end{tcbinfo}

\begin{tcbinfo}{GaN充電器——出差黨的好搭檔}
你可能注意到現在的快充頭越做越小，背後功臣是\textbf{GaN（氮化鎵）}這種新型半導體材料。相比傳統矽基充電器，GaN充電器有幾個明顯優勢：

\begin{itemize}
  \item \textbf{體積小}：同等功率下，GaN充電器可以做到傳統矽基的1/2甚至1/3大小，65W功率的GaN頭比蘋果老的5W頭還小。
  \item \textbf{效率高、發熱低}：氮化鎵的開關損耗遠低於矽，能量轉換效率可達95\%以上，充電時不容易燙手。
  \item \textbf{適合出差}：一個65W或100W的GaN多口充電器，能同時給手機、平板、筆記本充電，比帶三塊"磚頭"輕便得多。
  \item \textbf{選購建議}：認準支援PD協議的GaN充電器（65W起步可充筆電），知名品牌如倍思、安克、努比亞、小米等均可。注意核對資料線是否匹配（高功率需用E-Marker晶片的Type-C線）。
\end{itemize}

\end{tcbinfo}

\booksection{三、快充協議全解析：為什麼你的充電器充不快？}

你可能遇到過這種情況：明明手機支援快充，換了充電器卻變得很慢。這是因為快充不是簡單的“功率大就行”，而是需要充電頭、資料線、手機三方“說同一種語言”——這就是快充協議。

\booksubsection{快充的基本原理}

充電功率 = 電壓 × 電流（P = U × I）。要提升功率，要麼提高電壓，要麼提高電流，或者兩者都提高。不同協議選擇了不同的技術路線。

快充的工作流程：

\begin{enumerate}
  \item \textbf{握手階段}：手機和充電頭透過資料線通訊，交換各自支援的協議和功率資訊。
  \item \textbf{協商階段}：雙方確定一個共同支援的最高功率檔位。
  \item \textbf{充電階段}：充電頭按協商好的電壓電流輸出，手機端進行電源管理。
  \item \textbf{調整階段}：隨著電量增加，手機會逐步降低充電功率，保護電池。
\end{enumerate}

\booksubsection{主流快充協議詳解}

\textbf{USB PD（Power Delivery）}——通用協議之王

\begin{itemize}
  \item \textbf{特點}：USB-IF組織制定的官方標準，覆蓋範圍最廣。PD 3.0支援最高100W（20V/5A），PD 3.1擴充到240W（48V/5A）。
  \item \textbf{電壓檔位}：5V、9V、15V、20V，PD 3.1還支援28V、36V、48V。
  \item \textbf{相容性}：iPhone全系、谷歌Pixel、部分安卓旗艦、膝上型電腦、Switch遊戲機都支援PD協議。
  \item \textbf{優勢}：通用性強，一個充電頭可以給多種裝置充電；支援PPS（可程式設計電源），可以精細調節電壓電流。
  \item \textbf{劣勢}：需要Type-C介面和專用晶片，成本較高。
\end{itemize}

\textbf{高通QC（Quick Charge）}——安卓陣營的元老

\begin{itemize}
  \item \textbf{QC 3.0}：支援3.6V-20V以0.2V為步進的可調電壓，最高18W。
  \item \textbf{QC 4.0/4+}：相容USB PD協議，支援最高27W，加入雙路充電技術。
  \item \textbf{QC 5.0}：支援最高100W+，0-50\%僅需5分鐘，大幅改善發熱。
  \item \textbf{相容性}：使用高通驍龍處理器的手機大多支援，但各品牌通常會優先使用自己的私有協議。
\end{itemize}

\textbf{OPPO VOOC/一加WARP/真我Dart}——低壓大電流的代表

\begin{itemize}
  \item \textbf{技術路線}：低壓大電流（5V/4A起步），發熱更小，效率更高。
  \item \textbf{VOOC 1.0}：20W（5V/4A），2014年釋出，開創快充時代。
  \item \textbf{SuperVOOC}：50W-65W，採用雙電芯串聯設計。
  \item \textbf{SuperVOOC 2.0}：65W，30分鐘充滿4000mAh電池。
  \item \textbf{150W/240W超級閃充}：目前最快的量產快充技術。
  \item \textbf{特點}：需要專用充電頭和資料線（有額外觸點），不相容其他協議。
\end{itemize}

\textbf{華為SuperCharge}——高壓快充的演進

\begin{itemize}
  \item \textbf{FCP（Fast Charge Protocol）}：9V/2A，18W，相容QC 2.0。
  \item \textbf{SCP（Super Charge Protocol）}：支援更高功率，從22.5W到66W再到100W。
  \item \textbf{技術特點}：採用電荷泵技術，將高壓轉換為低壓大電流輸入電池，效率高達98\%。
  \item \textbf{相容性}：華為/榮耀手機專用，部分充電器支援PD協議。
\end{itemize}

\textbf{小米澎湃快充}——自研協議的崛起

\begin{itemize}
  \item \textbf{澎湃P1晶片}：小米自研的快充管理晶片，支援120W有線快充。
  \item \textbf{澎湃G1晶片}：電池管理晶片，與P1配合實現更高效的快充。
  \item \textbf{技術特點}：採用雙電芯串聯+電荷泵方案，支援多種協議相容。
\end{itemize}

\textbf{蘋果快充}——保守但穩定

\begin{itemize}
  \item \textbf{協議}：僅支援USB PD協議。
  \item \textbf{功率}：iPhone 8及以後支援快充，最高約27W（iPhone 15 Pro Max可達27W+）。
  \item \textbf{特點}：功率相對保守，但充電曲線最佳化好，發熱控制出色。
  \item \textbf{配件要求}：需要USB-C to Lightning線（或USB-C to USB-C線）和PD充電頭。
\end{itemize}

\booksubsection{快充協議對比表}

\begin{center}
\begin{tabularx}{\linewidth}{|X|X|X|X|}
\hline
\textbf{協議} & \textbf{最高功率} & \textbf{技術路線} & \textbf{相容性} \tabularnewline
\hline
USB PD 3.1 & 240W & 高壓/多檔位 & 最廣泛 \tabularnewline
QC 5.0 & 100W+ & 高壓可調 & 高通機型 \tabularnewline
VOOC/Dart & 240W & 低壓大電流 & OPPO系專用 \tabularnewline
SuperCharge & 100W & 電荷泵 & 華為/榮耀 \tabularnewline
蘋果快充 & \textasciitilde{}27W & PD協議 & iPhone/iPad \tabularnewline
\hline
\end{tabularx}
\end{center}

\booksubsection{無線充電協議}

無線充電也有自己的協議標準：

\begin{itemize}
  \item \textbf{Qi標準}：WPC（無線充電聯盟）制定，最通用的標準。基礎Qi支援5W-15W。
  \item \textbf{Qi2標準}：2023年釋出，引入磁吸定位（類似蘋果MagSafe），支援15W+。
  \item \textbf{私有協議}：各品牌有自己的高速無線充電，如小米50W無線充、華為50W無線充等，需要專用充電器。
\end{itemize}

\booksubsection{快充選購指南}

\begin{itemize}
  \item \textbf{看協議匹配}：購買充電器時，確認它支援你手機的快充協議。第三方充電器通常會標註支援的協議。
  \item \textbf{看資料線}：大功率快充需要專用資料線。普通資料線可能只支援60W以下。
  \item \textbf{看認證}：蘋果使用者建議選擇MFi認證的資料線；安卓使用者選擇品牌原裝或知名第三方。
  \item \textbf{看發熱}：快充時手機發熱是正常的，但如果燙手就不正常了，建議暫停充電或更換充電器。
\end{itemize}

\begin{tcbwarn}{快充安全提示}
\begin{itemize}
  \item 避免邊充邊玩大型遊戲，會加劇發熱。
  \item 不要使用劣質充電器和資料線，可能導致安全隱患。
  \item 充電時取掉厚重的手機殼，幫助散熱。
  \item 如果充電器或手機異常發熱、有異味，立即停止使用。
\end{itemize}

\end{tcbwarn}

\booksection{四、電池新技術：矽碳負極電池}

如果你留意2024-2026年的旗艦手機釋出會，“矽碳負極”這個詞幾乎成了高階電池的代名詞。這是鋰電池負極材料的一次重要升級，原理其實不復雜：

\begin{itemize}
  \item \textbf{傳統石墨負極}：手機電池一直用的是石墨做負極，技術成熟、穩定，但能量密度已接近理論上限（約372mAh/g）。
  \item \textbf{摻入矽材料}：矽的理論比容量是石墨的10倍左右（約4200mAh/g）。把少量矽摻入石墨裡做成"矽碳複合負極"，就能顯著提升整塊電池的能量密度。
  \item \textbf{能量密度提升}：同等體積下，矽碳負極電池的容量可以提升約\textbf{15\%-20\%}。這正是2024-2026年旗艦機能在7-9mm機身裡塞進5500-6000mAh大電池的關鍵原因。
\end{itemize}

\begin{tcbinfo}{哪些機型用上了？}
\begin{itemize}
  \item \textbf{榮耀}：Magic系列率先商用矽碳負極，Magic6 Pro、Magic7系列、Magic V3摺疊屏均採用。
  \item \textbf{小米}：小米14 Ultra、小米15系列搭載新一代矽碳負極電池。
  \item \textbf{vivo}：X100 Ultra、X200 Pro系列使用高密度矽碳負極方案。
  \item \textbf{華為/OPPO等}：2025-2026年旗艦普遍跟進，矽碳負極正在從"黑科技"變成高階機的標配。
\end{itemize}

\end{tcbinfo}

\textbf{對使用者的實際影響}：矽碳負極帶來的最大好處是“同體積電池容量更大”——手機可以做得更薄，同時續航不縮水；或者機身不變厚，但電池從4500mAh漲到5500mAh。換句話說，你不用為續航付“厚重”的代價了。需要注意的是，矽碳負極電池在低溫和滿充管理上仍需配合快充晶片做最佳化，這也是廠商建議長期保持80\%充電的原因之一。

\booksection{五、這才是真正的省電技巧}

\begin{itemize}
  \item \textbf{螢幕（耗電王）}：調低亮度、縮短自動鎖屏時間。OLED屏手機開深色模式能省不少電。
  \item \textbf{訊號差}：在電梯、地下室，手機會拼命搜訊號，耗電極快。
  \item \textbf{5G}：如果訊號不好，手動切回4G，續航會好很多。
  \item \textbf{定位服務}：非導航類App，定位許可權設為"僅使用時允許"。
  \item \textbf{限制後臺活動}：iPhone在設定 > 通用 > 後臺App重新整理；安卓在應用電池設定裡。
\end{itemize}

\begin{tcbpractice}{來檢查一下：給你的電池做個體檢}
\begin{enumerate}
  \item 檢視健康度：找到並記下你的電池最大容量。
  \item 分析耗電詳情：看電池設定裡的圖表和App列表，找出耗電元兇。
  \item 制定對策：限制後臺耗電異常的App。
\end{enumerate}

\end{tcbpractice}
